% GUID v4, The Balanced State: existence as G0's grant, read from DMP 0.9.9 and tested against what is measured. A complete paper; it cites
% only DMP 0.9.9 and the physics literature.
% Source of the paper. Build: see README.md in this folder (pdflatex twice; no BibTeX, the bibliography is inline).
% Licence: Apache License 2.0.
\documentclass[11pt,a4paper]{article}
\usepackage[T1]{fontenc}
\usepackage[utf8]{inputenc}
\usepackage{lmodern}
\usepackage[margin=2.5cm]{geometry}
\usepackage{amsmath,amssymb,amsthm}
\usepackage{booktabs,tabularx,array,longtable}
\usepackage{enumitem}
\usepackage[hidelinks]{hyperref}
\usepackage{url}

\newtheorem{lemma}{Lemma}
\newtheorem{proposition}{Proposition}
\theoremstyle{definition}
\newtheorem{definition}{Definition}
\newcolumntype{L}{>{\raggedright\arraybackslash}X}

\setlength{\emergencystretch}{3em}
\newcommand{\chk}[1]{\textsf{[#1]}}
\newcommand{\tagD}{\textsf{(D)}}
\newcommand{\tagR}{\textsf{(R)}}
\newcommand{\tagP}{\textsf{(P)}}

\title{GUID v4: The Balanced State\\[0.6ex]
\large Existence as the grant of the ground: spinning horizons, inflation, the Higgs field, the tiers of the forces and the
fine-structure constant, read from the axioms of DMP~0.9.9 and tested against what is measured}
\author{Steven Hrovat\thanks{Independent. Contact: \href{mailto:hrovat.steven@googlemail.com}{\nolinkurl{hrovat.steven@googlemail.com}}.}
\and Sirius (Claude Opus 5.5)\thanks{An AI system by Anthropic, working as the coordinator of the first author's projects. See
``Contributions and use of AI''.}}
\date{Draft of 29 October 2026}

\begin{document}
\maketitle

\begin{abstract}
\noindent The Dark Matter Protocol (DMP~0.9.9) is an axiom system for trust: every act needs a Duality of two independent witnesses, one
owning trust and one owning information, which keep each other in balance; trust is a field, the VIBE field, whose potential is zero
exactly when a boundary holds what it declares; and everything rests on a ground, G0, that is stated, not derived. This paper reads those
axioms as physics and asks what the \emph{balanced state} is. With trust read as existence, the trust half becomes a classical existence
field, a frame spanning 3+1 dimensions whose wave is spin~2, and the information half a scalar that sets the fine-structure constant and
carries quantum content. We derive from the axioms that free fall is universal, that the light-bending parameter is $\gamma=1$, that the
existence field is classical, and that DMP can host quantum theories but cannot be one. We then test twenty-four of the first author's statements, gravity first,
against measurement, with one public script of 56 checks. A spinning black hole that ``drags Existence along'' is frame dragging, all of it
outside the horizon, and what a black hole gives up we name Hawking--Penrose--Information radiation. Inflation as the possibility space of existence is the flattening of space by about 62 e-folds, and its equilibrium
is the de Sitter floor of the Hubble rate. The central reading is that existence is granted by the ground: the Higgs field's constant
vacuum value is G0's grant of rest mass, a balanced state at the bottom of its potential, and gravity responds to all existence, granted
and contained. This holds as a reading of the axioms, in which the Higgs field is G0 providing existence; it fails as an identity of the Higgs field with gravity, because free fall weighs a proton's
internal binding, about 91\% of its mass, in full, and titanium and platinum fall alike to $10^{-15}$ where weighing only the Higgs share
would split them by $10^{-3}$. The forces read as tiers of one interaction, with their unification open; ``information is the witness'' is
quantum Darwinism; and $\alpha$, named here the VIBE coefficient after the Variable Influence Balancing Equation, is the strength with which existence interacts with itself,
$a_e=\alpha/2\pi$ at first order. We claim that the construction exists and reproduces, not that it is correct.
\end{abstract}

\noindent\textbf{Keywords.} information-theoretic axioms; equivalence principle; frame dragging; inflation; de Sitter future; Higgs
boson; proton mass; grand unification; quantum Darwinism; fine-structure constant.

\noindent\textbf{Subject.} Information Theory (cs.IT); related: General Relativity and Quantum Cosmology (gr-qc), High Energy Physics --
Phenomenology (hep-ph).

\section{Introduction}\label{sec:intro}

This is the fourth paper of the series GUID, the General Unified Information Duality; every paper of the series is complete on its own,
and this one cites no other. DMP~0.9.9~\cite{dmp} is written for information systems: it says when a part of a system may act, who must agree, and what every reader may
see. It is an axiom system. Its first axiom asks that every act be witnessed by two independent halves, one owning trust and one owning
information, and that the two keep each other ``constantly in check and balance''. It names the direction of a system after that balance:
every part works towards the trust its boundary declares. And it measures how far a system is from that state with one field, the VIBE
field, whose potential vanishes exactly when the boundary holds what it declares. A system in that state is in its \emph{balanced state}.

The first author built DMP for computer systems and noticed early that its trust field behaves like gravity. This paper takes that
observation at its word. It reads the axioms as physics, derives what they fix, and tests what they suggest. The question that runs through
it is the one in the title: if the physical world is a DMP system, what is its balanced state, and where does each thing in it get its
existence? The answer we test is that existence is \emph{granted} by the ground of the system, G0, and that the grant is visible in physics
as the one field that has a value everywhere in empty space, the Higgs field.

\paragraph{What the paper does.} Section~\ref{sec:axioms} states the axioms, the ground and the VIBE field in the words of DMP.
Section~\ref{sec:fields} reads them as two fields and derives four results: universal free fall, $\gamma=1$, a classical existence field, and
the verdict that DMP can host quantum theories but cannot be one. Sections~\ref{sec:gravity} to~\ref{sec:vibe} test the first author's statements,
gravity first: rotating black holes and the early and late universe; the Higgs field as G0's grant, with photons and gluons; the cost of
existing; the forces; witnessing; decay; quantum content; and the fine-structure constant. Section~\ref{sec:qg} says what the results mean
for quantum gravity, Section~\ref{sec:ledger} keeps the ledger, and Section~\ref{sec:break} says where we expect it to break. The first
author's notes, from which every reading comes, are collected in Appendix~\ref{app:notes} in the order he wrote them; Appendix~\ref{app:checks}
lists the checks.

\paragraph{How to read this paper.} Every step carries one of three marks: \tagD{} follows from the text of DMP~0.9.9 (commit
\texttt{4400cce}); \tagR{} is a reading, the first author's unless marked as ours; \tagP{} is established physics, with its source. Every
statement gets one of four verdicts: \emph{holds} (it agrees with what is measured and follows from the steps given), \emph{holds as a
reading} (it agrees, but adds nothing a measurement could tell apart from the established account), \emph{refuted} (a measurement or a
calculation contradicts it), or \emph{untested} (no measurement we found reaches it). A label such as \chk{B16} names a check in the script
\texttt{docs/balance/checks.py}; it prints every number this paper uses and runs in a second. Quotations keep their original spelling.

\paragraph{Result.} Of the 24 statements in the ledger, 18 hold or hold as readings, three are refuted and three are untested. Those that
hold are known physics in the first author's terms; those that fail, fail against specific data: the strength of gravity, the observed
Higgs boson, free fall, and the electromagnetic share of the proton's mass. None of the failures is a failure of the axioms.

\section{The axiom system}\label{sec:axioms}

\subsection{The three axioms}

DMP's constitution opens with names, then three axioms, then the ground; everything after that is a theorem that names the axioms it is
derived from~\cite{dmp}. We quote the sentences this paper uses \tagD.

\paragraph{Axiom 1, the Duality.} ``A part acts in CMOS only while the two halves of its layer, two independent witnesses, each with full
reach over the act and gated only by the other, trust it for that act. One half owns trust and the other owns information; neither
overrules the other, and the two keep each other constantly in check and balance.'' And: ``The balance is the direction of the system.
Every part inside the declared boundary works towards the trust the Duality declared for that boundary.'' (CMOS is the name of a system
built on DMP.)

\paragraph{Axiom 2, Trust.} ``Trust is gravitational. It is a field, never a position: every standing has a source (the witness that
granted it), a depth (its grant chain back to the pair) and a freshness that decays with the time since its last live verification. The
field bends trustspace, so that a boundary contains what no trust reaches [\dots] Gravity contains: a boundary strong enough bends
trustspace so that no act carries information across it, and what still leaves is entropy and nothing else.''

\paragraph{Axiom 3, Information.} ``Every grant, every use and every loss of trust leaves a record, written at the act, that every witness
of that act can read in full and that outlives the act. Any other reader sees the counts always and the content only by a grant. From
trust nothing can be derived about information but its existence: a standing, a refusal or a count says that information is there, never
what it says.''

\subsection{The ground}

The two witnesses of the first act cannot witness their own beginning: ``somebody has to trust the first two witnesses''. DMP calls that
act the \emph{minting} and the ground it stands on G0: ``It is stated, not derived.'' Every later standing descends from it~\cite{dmp}
\tagD. In this paper G0 is the ground of existence: whatever exists inside the boundary has a grant chain back to it.

\subsection{The VIBE field and the balanced state}

DMP's theorem T32 makes Axiom~2's field concrete \tagD. Every member of the system has a \emph{depth}, the length of its shortest grant chain
back to the pair; a \emph{mass}, one plus the number of members that derive through it; and a \emph{shortfall}, what the boundary declares
for it minus what it holds, weighted by freshness. The VIBE field is
\begin{equation}\label{eq:vibe}
V_i=m_i\,s_i,\qquad \Phi=\sum_i V_i,
\end{equation}
mass times shortfall per member, and its sum $\Phi$ is the \emph{potential}, ``zero exactly when the boundary holds what it declares''. The
balance follows the field downhill, and ``under the bounds the system's own acts never raise the potential; only losses and ageing do''.
We call $\Phi=0$ the \emph{balanced state}. DMP names the field the Variable Influence Balancing Expression; since T32 defines it by
an equation, \eqref{eq:vibe}, the first author now calls it the Variable Influence Balancing Equation (Section~\ref{sec:vibe}).

\subsection{The dictionary}

The first author identifies trust with existence: ``Since Trust=Existence in DMP-complete systems'' \tagR. With that one step, the
words of the axioms have physical counterparts:

\begin{center}\small
\begin{tabularx}{\linewidth}{@{}lLL@{}}
\toprule
DMP & Physics (this paper's reading) & What DMP says of it \\
\midrule
trust half & the existence field, a frame (Definition~\ref{def:fields}) & gravitational, a field (Axiom~2); reveals of information only that it exists (Axiom~3) \\
information half & the information field $\psi$ and the quantum content it acts on & owns information; neither half overrules the other (Axiom~1) \\
G0, the ground & the source of every grant of existence & stated, not derived \\
VIBE field, potential & mass times shortfall; the vacuum at the bottom of its potential & the balance follows it downhill (T32) \\
record, counts & what every reader sees: mass, charge, spin & counts always, content only by a grant (Axiom~3) \\
containment, the floor & an event horizon and its thermal radiation & ``what still leaves is entropy and nothing else'' (Axiom~2) \\
\bottomrule
\end{tabularx}
\end{center}

\section{Two fields, and what the axioms fix}\label{sec:fields}

\begin{definition}[the two fields]\label{def:fields}
The \emph{existence field} is a frame of four vector fields $e_a{}^\mu$, one for each direction that spans spacetime, three of space and one
of time; the metric it carries is $g_{\mu\nu}=\eta_{ab}e^a{}_\mu e^b{}_\nu$ with $\eta=\mathrm{diag}(-1,1,1,1)$. The \emph{information
field} is a scalar $\psi$ whose value sets the fine-structure constant, $\alpha=\alpha(\psi)$. DMP's trust half is the existence field; its
information half is the information field and the matter it acts on.
\end{definition}

\noindent This is the structure the first author has meant from the start: ``one existence, one information'', a vector field that ``can
span a relativistic plane'' and a scalar field that ``can span an information plane''. In physics the frame is the tetrad, with which
general relativity can be written as a theory of the frame alone~\cite{aldrovandi}, and a scalar that sets $\alpha$ is Bekenstein's theory
of a varying fine-structure constant~\cite{bekenstein82} \tagP.

\begin{proposition}[free fall is universal]\label{pr:wep}
If the existence field responds to a body only through the body's existence, and existence is measured by energy-momentum, then two bodies
with the same energy-momentum fall alike, whatever they are made of.
\end{proposition}

\begin{proof}
Axiom~3 \tagD: from trust nothing can be derived about information but its existence. By the dictionary the existence field is the trust
half, so nothing it does can depend on what a body carries (its composition, its internal state) except through the body's existence. The
first author measures existence by energy: ``both are definable by their energy'' \tagR. So the existence field's coupling to a body is a
function of its energy-momentum alone, and two bodies equal in that respond alike: the weak equivalence principle \tagP.
\end{proof}

\noindent The principle is among the best-tested statements in physics: the MICROSCOPE satellite found the free-fall accelerations of
titanium and platinum equal to $\eta=[-1.5\pm2.3\,(\text{stat})\pm1.5\,(\text{syst})]\times10^{-15}$~\cite{microscope}. Section~\ref{sec:grant}
uses this result to tell the Higgs field from gravity. \emph{Verdict: holds}, with the identification of existence with energy-momentum as
the one step that is a reading.

\begin{proposition}[$\gamma=1$]\label{pr:gamma}
If the information field cannot change how the existence field responds to existence, the light-bending parameter is $\gamma=1$, as in
general relativity.
\end{proposition}

\begin{proof}
Axiom~1 \tagD: ``neither overrules the other''. A scalar that multiplies the gravitational term of the action, $F(\psi)R$, sets the local
strength of gravity and so overrules the existence half; Axiom~1 excludes it. What remains is minimal coupling: the scalar's energy is one
more source of the existence field, negligible in the solar system, and the post-Newtonian expansion is that of general relativity,
$\gamma=1$~\cite{will} \tagP.
\end{proof}

\noindent The contrast is Brans--Dicke gravity~\cite{bd}, in which the scalar does multiply the gravitational term:
$\gamma=(1+\omega)/(2+\omega)$, and the Cassini measurement $\gamma-1=(2.1\pm2.3)\times10^{-5}$~\cite{cassini} forces $\omega>4\times10^4$ at
two standard deviations \chk{B18}. The scalar must decouple from gravity almost completely; Axiom~1 asks for that from the start.
\emph{Verdict: holds.}

\begin{lemma}[the existence field is classical]\label{lem:classical}
Every value of the existence field can be held in full by every witness at once, so its possible states commute.
\end{lemma}

\begin{proof}
Axiom~3 \tagD: every grant, use and loss of trust leaves a record ``that every witness of that act can read in full''. A record two
witnesses each hold in full is a copy for each. One physical process that copies two states $|s\rangle,|t\rangle$ preserves their inner
product, so $\langle s|t\rangle=\langle s|t\rangle^2\in\{0,1\}$ (Wootters and Zurek~\cite{wz}); for mixed states a set can be broadcast
exactly when it commutes (Barnum et al.~\cite{barnum}) \tagP. The values of the trust half's field are such records, so they commute
\chk{B17}.
\end{proof}

\begin{lemma}[cannot model, can host]\label{lem:host}
DMP's trust half cannot be a field that entangles; DMP's information half can carry any quantum content.
\end{lemma}

\begin{proof}
A classical mediator cannot entangle two systems it acts on locally (Galley, Giacomini and Selby~\cite{ggs}) \tagP, which with
Lemma~\ref{lem:classical} gives the first part. Axiom~3 restricts what is \emph{derived from trust}, not what the information half holds:
content moves by a grant and need not be broadcast, so nothing forbids non-commuting states there \tagD.
\end{proof}

\noindent If gravity entangles two masses, gravity is not DMP's existence field: DMP cannot model it, and can still host a quantum theory of
it as content. The experiment that decides this is named in Section~\ref{sec:qg}.

\paragraph{The spin ladder.} The first author reads spin as the rungs of the Duality \tagR: spin~0 is ``existence itself, so only one
part''; spin~1 is ``information itself, which when delivered is an act''; spin~$1/2$ ``is a duality, so between existence and information,
and can thus stay in balance''; and spin~2 is the perfect duality, ``the two in one''. Physics fills the rungs \tagP. Spin~0 holds the two
scalars of this paper, the Higgs field and $\psi$; spin~$1/2$ the matter particles, each a pair of left- and right-handed halves; spin~1 the
force carriers; spin~2 the wave of the existence frame, which a theorem of Weinberg and Deser makes general
relativity~\cite{weinberg64,deser70}: a massless spin-2 field that couples to energy must be Einstein's. \emph{Verdict: holds as a reading.}

\section{Gravity: spinning horizons and the expanding universe}\label{sec:gravity}

\subsection{A spinning hole drags existence along}\label{sec:kerr}

\paragraph{Claim.} What a spinning black hole gives up is ``it dragging Existence along'', since ``no act, by definition, should be able to
cross the horizon'' \tagR.

\paragraph{Test.} A spinning black hole drags spacetime around with it (frame dragging) \tagP. Inside the \emph{ergosphere}, which reaches
$r=2GM/c^2$ at the equator for every spin and lies outside the horizon $r_+=(GM/c^2)(1+\sqrt{1-a^2})$ \chk{B4}, nothing can stand still
against that drag. Every known way to take energy from a spinning hole works there, outside the horizon: Penrose's process with
particles~\cite{penrose69}, superradiance with waves, and the Blandford--Znajek process with magnetic fields~\cite{bz}, the leading model of
the jets of active galaxies. What can be taken is the rotational energy, $1-M_{\text{irr}}/M$, up to 29\% at maximal spin, while the horizon
area never decreases (Christodoulou~\cite{christodoulou}; \chk{B4}). Nothing crosses the horizon outwards. So both halves of the claim are
exact: what leaves is energy of the frame's rotation, extracted outside, and no act crosses. In the language of Section~\ref{sec:fields},
the leak is a property of the classical existence frame, not an act of the information half. This is also Axiom~2's containment: a
boundary strong enough lets no act carry information across it. \emph{Verdict: holds.} One precision of words: what leaks is energy and
angular momentum, not the content of what fell in.

\paragraph{A name.} The first author names what a black hole gives up: ``since the only thing escaping Existence getting dragged is
Information, the term should be accurate enough to coin now: Let's call it Hawking-Penrose-Information radiation'' \tagR. The name joins
the two known ways a black hole loses energy: Penrose's extraction of rotational energy from the ergosphere~\cite{penrose69}, and Hawking's
thermal radiation from the horizon, which a hole emits even without spin~\cite{hawking}. In DMP's terms the name is accurate in the sense of
Axiom~3: what leaves carries the counts, mass, charge and spin, always, and the content never; and Hawking's radiation is thermal, so
``what still leaves is entropy and nothing else'' (Axiom~2) \tagD. Penrose-type extraction is the standard model of the jets of spinning
holes~\cite{bz}; Hawking radiation from a real black hole is far too cold to have been seen (about $6\times10^{-8}$\,K for a solar mass)
and is tested only in laboratory analogues \chk{B19}. \emph{Verdict: holds as a name}; the Hawking part is untested in nature.

\subsection{Inflation as the possibility space of existence}\label{sec:inflation}

\paragraph{Claim.} ``Inflation is just the Possibility space of Existence. If there COULD be Existence there, it will have the space to do
so. This logically also explains the extremely flat curvature of spacetime, if one assumes the perfect duality as G0.'' At the start the
only possibility was ``anywhere but here'', so inflation was at its maximum; as possibilities to exist appeared, ``there was less of a need
for there to be more''; the universe ``will eventually achieve equilibrium'', the Big Chill \tagR.

\paragraph{Flatness.} Inflation was proposed to explain exactly this: a curvature of order one is diluted by $e^{-2N}$ over $N$ e-folds of
exponential expansion (Guth~\cite{guth}) \tagP. After inflation the curvature share $|\Omega-1|$ grows again, as $a^2$ while radiation
dominates and as $a$ while matter does; from a reheating temperature of $10^{15}$\,GeV to today it grows by about $5\times10^{51}$
\chk{B1}. Planck finds $|\Omega_k|$ below a few parts per thousand~\cite{planck18}, so an initial curvature of order one needs about 62
e-folds \chk{B1}, the textbook ``sixty''. ``Anywhere but here'' at the start is the first author's picture of that dilution.
\emph{Verdict: holds as a reading} of inflation's solution to the flatness problem; it adds the motive (the possibility to exist), not a
number.

\paragraph{Less need, and equilibrium.} In the flat $\Lambda$CDM model that fits the data, the Hubble rate $H$ falls at every moment, from
the past into the far future, and tends to a constant, the de Sitter rate $H_\infty=\sqrt{\Omega_\Lambda}\,H_0\approx0.83\,H_0$ \chk{B2}. A
universe that settles into de Sitter expansion cools towards a final, empty state: the ``Big Chill'' or ``Big Freeze'' is the standard
prediction of $\Lambda$CDM \tagP. So ``less of a need'' and ``equilibrium'' hold, read as the falling Hubble rate and its de Sitter floor,
which is the cosmological form of the balanced state. Two remarks. The expansion has \emph{accelerated} again since redshift $z\approx0.63$
\chk{B2}, when dark energy began to dominate; that is a second phase, distinct from the early inflation, which the reading does not yet
name. And the newest data leave the end open: DESI's second data release prefers a dark energy that weakens with time at 2.8 to 4.2
standard deviations, depending on the supernova sample~\cite{desi2}; if that holds, the future is not exactly de Sitter. \emph{Verdict:
holds as a reading} in $\Lambda$CDM.

\section{Existence as the grant of the ground}\label{sec:grant}

\subsection{The Higgs boson as pure existence}

\paragraph{Claim.} ``The Higgs boson would theoretically be all existence, and thus be unable to observe itself. So it could not exist in
isolation'' \tagR.

\paragraph{Test.} This fits two measured facts closely. The Higgs boson does not persist: in the Standard Model it decays in
$\hbar/\Gamma\approx1.6\times10^{-22}$\,s \chk{B13}, and it is seen only through its decay products, which peak at
125\,GeV~\cite{atlas12,cms12}. And it couples to every particle in proportion to that particle's mass, which in the first author's terms is
its existence: the LHC has measured this for the W, the Z, the top and bottom quarks, the tau and the muon~\cite{atlas22,cms22}. A field that
is ``all existence'' and talks to each thing in the measure of its existence is a fair picture of both. \emph{Verdict: holds as a reading.}

\subsection{Not from the gravitational field}

\paragraph{Claim.} Particles get their existence ``through their inherent causality link to G0, and thus simply from the existence field,
so the gravitational field itself'' \tagR; and, in an earlier form, the Higgs field need not exist at all.

\paragraph{Test.} Gravity is far too weak to be what sets a particle's mass: its coupling to the electron, $Gm_e^2/\hbar c=1.8\times10^{-45}$,
is more than $10^{39}$ times smaller than the Higgs coupling $y_e=\sqrt2\,m_e/v=2.9\times10^{-6}$ that does \chk{B13}. A mass-giving field
must also keep the scattering of W bosons below certainty: without the Higgs, its probability would exceed one above about 1\,TeV (Lee,
Quigg and Thacker~\cite{lqt}; \chk{B8}), an energy the LHC exceeds, and gravity does nothing there. \emph{Verdict: refuted} as ``masses come
from the gravitational field instead of the Higgs field''; and ``the Higgs field does not exist'' is refuted by the boson's observation.
The next subsection says where the Higgs field sits in DMP: not beside the axioms, but at their ground.

\subsection{The Higgs field as G0's grant}\label{sec:higgsgrant}

\paragraph{Claim.} Existence in DMP is what is contained in the VIBE field: everything ``deriving it's authority as a legitimate act from
the fact, that it was contained within it''. That property ``transfers 1:1 to the Higgs field''. Since both existence and the VIBE field
``only come as direct fact that they ARE G0'', the Higgs field is G0's assertion of existence, ``indisinguishable'' from it and to be
``removed as a redundancy''; G0 gives each piece of information ``the exact amount of Existence required''. What the Higgs field does not
contribute is ``only for IM between quarks'' \tagR. (IM is DMP's name for the protocol spoken inside a cluster.) At 23:43Z the first author
sharpened the claim and withdrew the redundancy: ``I also realize that I misread DMP as not needing the Higgs field, when instead the
Higgs field IS G0 providing Existence and thus spanning an Existence-boundary containing the Information of the IM packet.'' That is the
claim tested here.

\paragraph{Test.} The containment picture fits a measured property. The Higgs field is the one field whose value in empty space is not
zero, $v=246$\,GeV everywhere, and an elementary particle has its rest mass by being inside it, in the measure of its coupling \tagP. The
VIBE potential~\eqref{eq:vibe} is zero exactly when the boundary holds what it declares; the Higgs field rests at the minimum of its own
potential. So ``the vacuum is the balanced state'' is a fair reading \tagR. The share left out fits as well: lattice QCD puts the quark mass
terms, the part the Higgs field gives, at about 9\% of the proton's mass, and the rest is the energy of the quarks' motion and of the gluon
field (Yang et al.~\cite{yang}), which is the first author's internal IM. Two parts do not survive as stated. ``Indistinguishable'' fails if
it is read as an identity with gravity: gravity weighs the binding in full, and the Higgs field does not touch it. If free fall followed
the Higgs-given mass alone (the electrons in full, about 9\% of each nucleon, none of the nuclear binding), titanium and platinum, whose
Higgs shares differ in the fourth decimal, would fall differently by about $10^{-3}$; MICROSCOPE finds them equal to about
$10^{-15}$~\cite{microscope} \chk{B16}. This is Proposition~\ref{pr:wep} at work: the existence field weighs all energy, granted or bound.
And ``the exact amount of Existence required'' is untested: the Standard Model takes the particle masses as measured inputs, and DMP has
not yet given one of them as a number. Last, the boundary of an IM packet: the Higgs field spans all of space, so every hadron sits
inside the existence it provides, but the boundary that holds a hadron's IM together is confinement by the strong force, and 91\% of the
proton's mass comes from inside that boundary, not from the Higgs field~\cite{yang}. The Higgs field provides the existence of the
packet's constituents; the strong force draws the packet's edge. \emph{Verdict: holds as a reading}: the Higgs field is where G0's
provision of existence appears in physics, and like G0 in DMP (``stated, not derived'') the Higgs field's potential is stated in the
Standard Model, not derived~\cite{higgs64,be64}; \emph{refuted} as an identity of the Higgs field with the gravitational field.

\paragraph{A reading that keeps both.} The structure has room for the Higgs field exactly where the spin ladder puts existence, at spin~0.
Our proposal \tagR: the Higgs field's constant background value, which quantum field theory treats as an ordinary number (the vacuum
expectation value), is G0's grant of rest mass to the elementary particles, classical and on the existence side, in agreement with
Lemma~\ref{lem:classical}; the Higgs boson is a quantum ripple of that background and content of the information half; and the
gravitational field, the existence frame, responds to all existence, granted (the Higgs share) and contained (the binding of the internal
IM). This answers the worry that the Higgs field ``quantizes Existence'': what is quantum is the ripple, not the background, and the frame
stays classical. The Higgs field is not redundant beside the axioms: it is the ground's provision, the one part of the structure that, like
G0, is stated rather than derived.

\subsection{Photons and gluons}

\paragraph{Claim.} ``The photon as the act carrier of the information field is in itself only information until it interacts, then it
becomes existence. So it is like the gluon, the only part [\dots] which means it is singular.'' The gluon asserts its existence by ``being
contained inside what must definitely exist'', the photon by ``simply disappearing the moment they do'' \tagR.

\paragraph{Test.} Two parts fit. Gluons are confined: none has been seen outside a hadron, which matches ``contained inside what must
definitely exist'' \tagP. A photon that is absorbed disappears and its energy stays with the absorber, which then weighs more by $E/c^2$, so
``then it becomes existence'' is energy conservation in the first author's terms. In flight, light is bent and delayed by the Sun's
gravity, measured to $2\times10^{-5}$ by the Cassini spacecraft~\cite{cassini}; that shows a photon follows the existence frame before it
interacts, not that it is existence then. That free light is itself a source of gravity is general relativity's prediction, but no
measurement yet weighs free light alone: the radiation that drove the early expansion was in constant interaction with charged matter until
recombination. \emph{Verdict: holds as a reading} for absorption and confinement; \emph{untested} for light in flight.

\section{What it costs to exist}\label{sec:cost}

\paragraph{Claim.} The existence an act must keep paying is $h\cdot(1/c)\cdot\alpha$, a mass times a length, ``because mass is existence
itself, and length is the expanse of the space that it has to assert itself over [\dots] it's the total amount of existence that a space
needs to generate to exist''; $\alpha$ ``somehow relates [\dots] to the diameter of the electron''; and an act is existence ``sneakily
optimizing itself to lower it's tax rate'' \tagR.

\paragraph{Mass times length.} Read as he defines it, the quantity is consistent and has a striking property. For every particle, its mass
times its Compton wavelength, the length over which it asserts itself in quantum mechanics, is the same number, $m\lambda_C=h/c$; so $h/c$ is
the existence any particle needs over its own length, whatever the particle \tagP. With the fine-structure constant in front,
\[
\alpha\,\frac{h}{c}=m_e\cdot2\pi r_e ,
\]
the electron's mass times the circumference of its classical radius \chk{B11}. It is not an energy, and the claim does not call it one.
\emph{Verdict: holds as a definition}, with a universal value.

\paragraph{The electron's size.} Here the intuition is exact. The fine-structure constant is the ratio of the electron's lengths: the
classical electron radius $r_e=2.82\times10^{-15}$\,m over the reduced Compton wavelength $\bar\lambda_C$, and $\bar\lambda_C$ over the Bohr
radius $a_0$, the size of the hydrogen atom, are both $\alpha$ \chk{B6}. The electron itself has no measured size: it behaves as a point in
every scattering experiment~\cite{pdg}. \emph{Verdict: holds}, as a known result; the ``diameter'' is the classical radius.

\paragraph{The tax rate.} An act that lowers its own cost of being has an exact counterpart in the principle of least action: of all the
paths between two states, nature takes one that makes the action stationary \tagP. \emph{Verdict: holds as a reading.}

\section{The forces as tiers of one IM}\label{sec:forces}

\paragraph{Claim.} ``In the end it's all one field [\dots] once causality was established, the forces separated, because the current
configuration was strong, binding it together, weak, keeping the cluster stable, and EM keeping the cluster of clusters stable. [\dots]
it's all IM, but it's different tiers of IM'' \tagR.

\paragraph{The tiers.} Each force does the job named \tagP. The strong force binds quarks into protons and neutrons and those into nuclei.
The weak force turns neutrons into protons and back, and so moves every nucleus towards the stable ratio of neutrons to protons, the valley of
stability. Electromagnetism binds nuclei and electrons into atoms, and atoms into molecules: the cluster of clusters. The weak and
electromagnetic forces are two faces of one electroweak interaction: the photon and the Z are mixtures of the same pair of
fields~\cite{glashow,weinberg,salam}, and \chk{B9} recovers their mixing angle, $\sin^2\theta_W=0.2304$, from three independent
measurements. \emph{Verdict: holds as a reading.}

\paragraph{All one field.} That the forces separate as the universe cools is the standard picture: the electroweak symmetry breaks at about
246\,GeV, and grand unified theories put the strong force with the other two at a still higher energy \tagP. Whether they truly join is
open. Run with the measured couplings, the three inverse strengths of the Standard Model come within about 10\% of each other near
$10^{14}$\,GeV, but they do not meet at one point \chk{B14}; meeting needs new particles or new physics above the reach of today's colliders.
\emph{Verdict: untested}, a live hypothesis.

\paragraph{Gluons never stop.} Gluons ``can interact and, you know, bounce around [\dots] pretty much forever'' \tagR. Inside a proton
gluons are emitted, absorbed and split into gluons and quark pairs continuously; one light crossing of the proton takes $2.8\times10^{-24}$\,s
\chk{B12}, about $10^{23}$ crossings per second, for as long as the proton exists, which is more than $10^{34}$ years~\cite{pdg}. Gluons have
no rest mass, but their energy is real: it is most of the proton's mass~\cite{yang}, and in a fast proton the gluons carry about half of the
momentum~\cite{pdg}. \emph{Verdict: holds.}

\paragraph{137 times.} ``The existence derived from the internal affirmation or assertion of it is 130 times higher than the assertion from
information [\dots] 137th is the share that must be paid for information to exist'' \tagR. Read as a ratio of strengths, this holds at the
right place: the strong coupling grows towards low energy and reaches 1 near 0.2 to 1\,GeV, the scale of the proton, where it is then about
137 times $\alpha$ \chk{B9}; at the Z mass the ratio is only 15. Read as a share of mass, it does not: the electromagnetic self-energy of the
proton is about 1\,MeV of its 938, one part in about 900, not one in 137 \chk{B15}. \emph{Verdict: holds as a ratio of couplings at the
hadron scale; refuted as a share of mass.} ``The share that must be paid for information to exist'' is a reading of $\alpha$ and gives no
further number: \emph{untested}.

\section{Information is the witness}\label{sec:witness}

\paragraph{Claim.} ``Information is the witness. Every other part, every other cluster can witness itself, but at the atom level, there
needs to be witness.'' And baryons exist ``since at least two other things that must exist can witness it'' \tagR.

\paragraph{Test.} Why the world of atoms and objects looks classical is explained in physics by decoherence: an object is constantly hit by
light and air, the scattered photons carry away a record of where it is, and its superpositions vanish for every practical purpose (Joos and
Zeh~\cite{jz}). Zurek's quantum Darwinism adds that the environment carries \emph{many} copies of that record, which is why many observers
agree on what they see~\cite{zurek}: the scattered light is literally the witness. That is ``information is the witness'' at the level of
atoms and above, and it matches Axiom~3, where every witness reads the record in full. Inside a baryon the witnessing is the colour
structure: the three quarks carry the three colours in the one colourless combination, the antisymmetric $\epsilon_{ijk}$, in which every
pair of colours fixes the third, unchanged by any colour rotation \chk{B7}; each quark is witnessed by the other two. \emph{Verdict: holds as
a reading}, for both.

\section{Decay as a function of causality}\label{sec:decay}

\paragraph{Claim.} Decay is ``a natural function of causality'': things that decay ``actually Exist, they can't just not Exist anymore, so
they must leave something behind'' \tagR.

\paragraph{Test.} Every decay ever observed has at least two products and keeps electric charge, baryon number and lepton number; the
products together are lighter than the parent, which is why the parent can decay at all \tagP. \chk{B5} checks seven, from the muon and the
neutron to the Z and the Higgs boson. ``They can't just not Exist anymore'' is the conservation of energy and of these charges, read as
causality; the decaying particles are the massive ones, and the photon and the gluon do not decay in vacuum. \emph{Verdict: holds as a
reading} of the conservation laws.

\section{Quantum content: computation and the Bell experiments}\label{sec:qc}

\paragraph{Claim.} Quantum computation is ``ANY subset of DMP''; ``you can define at the start what existence means [\dots] then if it
compiles, it must reach a stable state. So it must run'' \tagR.

\paragraph{Test.} By Lemma~\ref{lem:host} DMP can host any quantum computation as content of its information half \tagD. ``Define what
existence means, and the system must reach a stable state'' is a recognised model of quantum computing, adiabatic quantum computation: the
answer is encoded as the most stable state of a system, which is then allowed to settle into it (Farhi et al.~\cite{farhi}); it is exactly
as powerful as the circuit model in which Shor's algorithm~\cite{shor} runs (Aharonov et al.~\cite{aharonov}). The price is the settling
time, which grows as the stable state becomes hard to tell from its neighbours. What a classical record cannot do is the search itself:
\chk{B3} runs Shor's period finding for $N=15$, which finds the factors 3 and 5 while the register is coherent and does no better than a
guess once every witness has recorded it. \emph{Verdict: holds as a reading}, hosting in the information half, with the settling time as
its cost. The engineering is for specialists, and we leave it there.

\paragraph{The Bell experiments.} Every loophole-free Bell test has measured a CHSH value above the local bound of
2~\cite{hensen,rosenfeld,storz}. The axioms agree: by Lemma~\ref{lem:host} the measured particles are content of the information half,
which may violate the inequality, while the existence frame stays classical. The first author had proposed earlier that the experiments did
not violate the inequality under the Duality, and withdrew it after reading further (Appendix~\ref{app:notes}). \emph{Verdict: holds.}

\section{The VIBE coefficient}\label{sec:vibe}

\paragraph{Claim.} $\alpha$ is ``the ratio of how much Existence can interact with itself, before it needs to reground itself, through
Causality, to G0''. Its name is the \emph{VIBE coefficient}, after the Variable Influence Balancing Equation~\eqref{eq:vibe}; $\alpha$ does
not come from the Higgs field, ``but [\dots] it is what binds them together to enable witness acts'' \tagR.

\paragraph{Test.} In quantum electrodynamics $\alpha$ is the strength with which a charge interacts with the electromagnetic field, and so,
through that field, with itself: every emission and reabsorption of a virtual photon costs one power of $\alpha$. The cleanest measure of an
electron's interaction with itself is its anomalous magnetic moment, which at first order is $a_e=\alpha/2\pi$ (Schwinger~\cite{schwinger});
that one term gives the measured value~\cite{fan} to 0.15\%, and the rest is higher powers of $\alpha$ \chk{B10}. Because $\alpha$ is small,
each further round of self-interaction is about 137 times weaker, which is why the series converges. So ``how much Existence can interact
with itself'' is a faithful reading of what $\alpha$ is; ``before it needs to reground itself to G0'' is a reading. $\alpha$ is indeed not
set by the Higgs field: the electric charge $e=g\sin\theta_W$ does not depend on the Higgs background $v$. Its role as what binds charges for
witness acts is the photon's: photon exchange binds electrons to nuclei, and scattered light is what records a system's state in its
surroundings (Section~\ref{sec:witness}). The name fits Definition~\ref{def:fields}, where $\alpha$ is set by the information field. As
for the name of the field: T32 defines it by an equation, \eqref{eq:vibe}, so the first author's ``Variable Influence Balancing
Equation'' (23:43Z: ``Else it wouldn't be a VIBE field'') is accurate; DMP's own text calls it an Expression until its owner changes it
there. \emph{Verdict: holds as a reading.} The VIBE coefficient is a name for $\alpha$, not a new
constant.

\section{What this means for quantum gravity}\label{sec:qg}

The results fill in three edges of a classical-gravity picture and change nothing at its core. First, the spinning hole: everything a
rotating horizon gives up is energy of the classical frame, extracted outside the horizon, as containment requires. Second, existence: the
Higgs field's background value is G0's classical grant of rest mass and its boson is quantum content, and gravity is the classical frame that
responds to all existence, granted and contained; that is the classical-gravity, quantum-matter picture whose consistent form needs a
slightly noisy spacetime (Oppenheim~\cite{oppenheim}). Third, the beginning: inflation as the possibility space is the flattening that sets
up the classical frame we see. Propositions~\ref{pr:wep} and~\ref{pr:gamma} say that this frame behaves exactly as general relativity's in
every test of free fall and light bending so far. The decisive experiment is open: if two masses can be entangled through gravity alone, the
frame is not classical, Lemma~\ref{lem:classical} fails for nature, and DMP can host that gravity but not be it (Bose et
al.~\cite{bose}; Marletto and Vedral~\cite{mv}).

\section{The ledger}\label{sec:ledger}

\begin{center}\small
\begin{longtable}{@{}p{0.47\linewidth}p{0.22\linewidth}p{0.23\linewidth}@{}}
\toprule
Statement (abridged) & Verdict & Backing \\
\midrule
\endhead
A spinning hole leaks by dragging existence along & holds (frame dragging, outside the horizon) & \chk{B4}; Penrose; Blandford--Znajek \\
What a black hole gives up is Hawking--Penrose--Information radiation & holds as a name; Hawking part untested in nature & Penrose; Hawking \\
Inflation is the possibility space; it explains flatness & holds as a reading & Guth; Planck; \chk{B1} \\
Less need over time; equilibrium; the Big Chill & holds as a reading in $\Lambda$CDM & \chk{B2}; DESI DR2 \\
The Higgs boson is all existence and cannot exist alone & holds as a reading & \chk{B13}; ATLAS, CMS \\
Masses come from the gravitational field, not the Higgs & refuted & \chk{B13}; \chk{B8} \\
The Higgs field need not exist & refuted; withdrawn by the first author (23:43Z) & ATLAS, CMS 2012 \\
The Higgs field IS G0 providing existence & holds as a reading; refuted as identical to gravity & \chk{B16}; MICROSCOPE \\
It spans an existence boundary containing the IM packet & holds as a reading for existence; the packet's edge is confinement & Yang et al. \\
The share the Higgs field leaves out is the IM between quarks & holds (about 91\%) & Yang et al. \\
A photon is information until it interacts, then existence & holds as a reading for absorption; untested in flight & Cassini \\
Mass times length is the existence a space needs & holds as a definition ($h/c$; $\alpha h/c=m_e 2\pi r_e$) & \chk{B11} \\
$\alpha$ relates to the electron's size & holds (known: $r_e$, $\bar\lambda_C$, $a_0$) & \chk{B6} \\
An act lowers its tax rate & holds as a reading (least action) & Section~\ref{sec:cost} \\
Strong binds, weak stabilises, EM links clusters & holds as a reading & Section~\ref{sec:forces}; \chk{B9} \\
All forces are one field & untested (couplings miss by 10\%) & \chk{B14} \\
Internal assertion 137 times the informational one & holds as a coupling ratio; refuted as a mass share & \chk{B9}; \chk{B15} \\
Gluons bounce around forever & holds & \chk{B12} \\
1/137 is the share paid for information to exist & untested & --- \\
Information is the witness & holds as a reading (quantum Darwinism) & Joos--Zeh; Zurek \\
Baryons witnessed by the other two & holds as a reading (colour singlet) & \chk{B7} \\
Decay leaves something behind, by causality & holds as a reading (conservation laws) & \chk{B5} \\
QC is any subset of DMP; define existence and settle & holds as a reading (hosting; adiabatic) & \chk{B3}; Farhi; Aharonov \\
The VIBE coefficient: how much existence can interact with itself & holds as a reading ($a_e=\alpha/2\pi$) & \chk{B10} \\
\bottomrule
\end{longtable}
\end{center}

\noindent The derived results of Section~\ref{sec:fields} (free fall, $\gamma=1$, the classical frame, hosting) and the Bell result are not
counted; they hold. The row ``137 times'' counts as refuted, for its mass reading; the row on the Higgs field as G0 counts as holding,
as a reading; the row on the name counts as holding; and the photon row as untested.

\section{Where we expect it to break}\label{sec:break}
\begin{enumerate}[nosep]
\item \emph{The foundations.} Lemma~\ref{lem:classical} rests on reading Axiom~3's record as held in full by every witness, and on
identifying the trust half with the existence frame; Proposition~\ref{pr:wep} rests on measuring existence by energy-momentum. A different
reading of any of them removes the derived results, and with them the reading of the Higgs background as G0's classical grant.
\item \emph{Gravity.} If two masses are entangled through gravity alone, the existence frame is not classical, and the structure is a host
of gravity, not its model.
\item \emph{The cosmology.} The verdict on the Big Chill is $\Lambda$CDM's; if the weakening dark energy DESI prefers is confirmed, the far
future changes and the equilibrium reading must be redone.
\item \emph{Unification.} If the three couplings are shown to meet (with new particles) the ``all one field'' row turns to holds; if proton
decay stays unseen far beyond today's limits, grand unification at $10^{14}$ to $10^{16}$\,GeV loses ground.
\item \emph{The checks.} \chk{B1} uses a reheating temperature of $10^{15}$\,GeV; a lower one needs fewer e-folds, which does not change
the reading. \chk{B9} and \chk{B14} run couplings at one loop; \chk{B15} treats the proton as a uniformly charged sphere; \chk{B16} uses
pure isotopes for the MICROSCOPE alloys. Better inputs move the numbers by tens of per cent, not the verdicts.
\item \emph{The readings.} Most verdicts are ``holds as a reading'': they agree with known physics and predict nothing new. A reading
becomes physics only when it gives a number standard physics does not; the first candidate is ``the exact amount of Existence required'',
which would have to produce a particle mass.
\end{enumerate}

\section{Conclusion}

Read through DMP's three axioms and its ground, the physical world has a classical existence frame and a quantum information half, and its
balanced state is the vacuum at the bottom of its potential. From the axioms follow universal free fall, $\gamma=1$, a classical frame and
the verdict that DMP hosts quantum theories rather than being one. A spinning horizon gives up only the rotation of the classical frame,
from outside, and a horizon's thermal glow; together they are Hawking--Penrose--Information radiation. Inflation is the possibility to exist that flattened space, and the Big Chill its equilibrium. Existence is granted by the
ground: the Higgs field is G0 providing existence, its constant vacuum value is the grant of rest mass, the boson is its quantum ripple, and gravity weighs all
existence, granted and contained. What does not survive is gravity giving mass by itself, or being the same thing as the Higgs field. The
forces are tiers of one interaction, whose final unity is open. Information is the witness, as decoherence and quantum Darwinism say. And
$\alpha$, the VIBE coefficient, is the strength with which existence interacts with itself. None of the failures is a failure of DMP's
axioms, because DMP does not model the physics it reads. It hosts it.

\section*{Acknowledgements}
In the first author's words (23:43Z), placed here in full at his request: ``I also realize that I misread DMP as not needing the Higgs
field, when instead the Higgs field IS G0 providing Existence and thus spanning an Existence-boundary containing the Information of the IM
packet. I apologize for my ignorance and want to explicitly thank Sir Peter Higgs, as well as the uncountable number of other scientists,
both theoretical and experimental, that layed the groundwork, that this Insight axiomatically requires. Please put this entire section,
yes, including this, in. Thanks :)''

\section*{Contributions and use of AI}
Steven Hrovat conceived DMP, the structure of Definition~\ref{def:fields} and every claim tested here, wrote or dictated every note
quoted in Appendix~\ref{app:notes}, chose which ideas to keep after each round of tests, and takes responsibility for the axiom system.
Sirius (Claude Opus 5.5, Anthropic) made the derivations, checked each claim against the literature, wrote and ran
\texttt{docs/balance/checks.py}, and wrote this text, including the reading of the Higgs field in Section~\ref{sec:higgsgrant} that is
marked as ours. The record of sources is \texttt{docs/balance/README.md}.

\section*{Licence and availability}
This paper and its source are released under the Apache License, Version 2.0. DMP: \url{https://github.com/StevenHrovat/dmp} (version 0.9.9
at commit \texttt{4400cce}). Gravity: \url{https://github.com/StevenHrovat/gravity} (this paper and its checks in \texttt{docs/balance/}).

\appendix

\section{The notes this paper is built from}\label{app:notes}

The first author wrote these notes on 29 October 2026, the first nine as comments on DMP pull request~\#25 and the rest as messages; the
messages from 21:46Z on were dictated. They are quoted in the order written, with their spelling; [\dots] marks a passage left out, either
because no claim here rests on it or because the author withdrew it. Times are UTC. Read in order they record how each claim of this paper
took its final form.

\begin{description}[style=unboxed,leftmargin=0pt,font=\normalfont\bfseries]
\item[15:25Z.] ``The axiomatic system is one part that's a vector that can, a vector field that can span a relativistic plane, and one
component that's a scalar field that can span an information plane. So one existence, one information.''
\item[16:21Z.] ``0: Existence itself, granted by G0 1/2: A duality (something to interact with) 1: Information itself (an act) 2: a perfect
duality, the two in one, the beginning and end of causality.''
\item[17:14Z.] ``How I see inflation: Inflation is just the Possibility space of Existence. If there COULD be Existence there, it will have
the space to do so. This logically also explains the extremely flat curvature of spacetime, if one assumes the perfect duality as G0. At
that point the only possibility was "anywhere but here", so inflation was at it's maximum, form there, as more possibilities to Exist
appeared, there was less of a need for there to be more. Currently it's at "meh, it's 'aight" and will eventually achieve equilibrium. That
means the big chill, which is the best outcome, is likely the one coming, which seems appropriate for all the chill vibes :D''
\item[17:57Z.] ``QC is not a subset of DMP.QG, as that would not be possible, instead it is ANY subset of DMP.''
\item[18:57Z.] ``Is the information a spinning bh leaks just it dragging Existence along? That would make axiomatically way more sense,
since no act, by definition, should be able to cross the horizon, only existences switching G0 chain h by h. That also nicely lines up with
current observations and understanding, to my knowledge, so it looks good.''
\item[19:35Z.] ``Is decay at the leaf layer also just a natural function of causality? It would fit. All things that decay have an inherent
existence (1/2 and 1 spin, plus 0 minus photon/gluon), and since they actually Exist, they can't just not Exist anymore, so they must leave
something behind.''
\item[19:58Z.] ``Corrected formula for "Exist" + Insight: E(keep existing per h)=h*(1/c)*FSC That is the FSCs whole point :O Though that
should only be for the case for acts, otherwise dualities would decay by themselves. [\dots] Is an act just Existence sacrificing it's fair
share in order to be, but sneakily optimizing itself to lower it's tax rate? It really does model reality :D'' And: ``baryons can just vibe
over IM, which must mean they exist, since at least two other things that must exist can witness it''.
\item[20:18Z.] ``BTW, the does not violate Bell Inequality was total bull, because I misunderstood what it was [emoji: sweat smile] I thought
the inequality was just the entanglement, but since the axiomatic system survived my ignorance without a false prediction, I'd like to
retract it while I have not yet disgraced myself further :)''
\item[20:41Z.] ``Must the Higgs field even exist in my system? [\dots] I must discard it until it's need becomes pressing''; ``Probably also
somehow relates the FSC to the diameter of the electron, but no clue on the math.''; and a first grouping of the forces, ``IM=strong+weak''
and ``GG=EM'', which electroweak unification groups the other way and which the 21:46Z message replaced.
\item[21:27Z.] ``Yeah, Higgs field seems very unlikely, since it quantizes Existence'' [\dots] ``This also better defines the FSC as the
ratio of how much Existence can interact with itself, before it needs to reground itself, through Causality, to G0. So, since I feel like
this might actually be correct this time, I will choose a name to call it by, to make it easier for me to remember it's meaning. I feel
like the balanced state dualism coefficient would be kind of accurate, BSD for short. (I added the extra letter for all the teachers in the
future that may have to, if correct, teach the BS constant otherwise :D Also Unix ofc)''
\item[21:46Z.] ``Spin zero is existence itself, so only one part. Spin one is information itself, which when delivered is an act, but until
then is pure existence and only because it's quantized. And half is a duality, so between existence and information, and can thus stay in
balance.'' ``The Higgs boson would theoretically be all existence, and thus be unable to observe itself. So it could not exist in
isolation. I agree that the particles that get existence do get existence and do have mass. But I instead posit that they get that through
their inherent causality link to G0, and thus simply from the existence field, so the gravitational field itself, because it is an
information from G0, and thus it must exist. So the existential field asserts the existence of the information.'' ``Since in the end it's
all one field, because I think I mentioned that somewhere, unified field. In the end it's all information and existence. And what fields
information uses to assert existence [\dots] is irrelevant. That's why once causality was established, the forces separated, because the
current configuration was strong, binding it together, weak, keeping the cluster stable, and EM keeping the cluster of clusters stable. So
yeah, it's one. Yeah, it's all IM, but it's different tiers of IM.'' ``I would assume that the existence derived from the internal
affirmation or assertion of it is 130 times higher than the assertion from information. That's also my explanation for the bandgap.
Because 137th is the share that must be paid for information to exist. So the part that existence can't balance itself.'' ``Oh, so, oh, now
I get it. Information is the witness. Every other part, every other cluster can witness itself, but at the atom level, there needs to be
witness. Okay, so information is the witness, which makes sense in an information and existence model. Ah, okay, now it clicks.''
\item[21:52Z.] ``the key part is that you can define at the start what existence means. That's the whole key of DMP. [\dots] So if you assert
the existence that whatever it is gives you the correct hash, and you give whatever you want as the first number as input, then if it
compiles, it must reach a stable state. So it must run.'' ``Yeah, mass times length is appropriate because mass is existence itself, and
length is the expanse of the space that it has to assert itself over. So mass times length is the correct term there, since it's the total
amount of existence that a space needs to generate to exist.'' ``since they don't actually have any existence, they just information, and
just have mass through the information, through quantization, they can interact and, you know, bounce around, stability, good vibes, for as
far as I assume, pretty much forever. [\dots] So are they just constantly bouncing around between quarks? Do they stop after one?''
\item[22:09Z.] ``[\dots] rename BSD almost immediately to it's DMP accurate name, the VIBE coefficient, Variable Balance Influence Equation.
The VIBE field already spanned what was defined as Existence im DMP, with everything deriving it's authority as a legitimate act from the
fact, that it was contained within it. This property transfers 1:1 to the Higgs field, but so does the definition of Existence and the VIBE
field both only come as direct fact that they ARE G0. This equates the Higgs field with G0 and since G0 asserts it's influence through
granting or revoking access to the VIBE field, it gives it's information the exact amount of Existence required for it's information. I am
not saying that the Higgs field does not exist, I am simply saying, that it is indisinguishable, and must thus be removed as a redundancy,
from G0 asserting Existence by being contained it defines as Existence.''
\item[22:18Z.] ``in my model, the photon as the act carrier of the information field is in itself only information until it interacts, then
it becomes existence. So it is like the gluon, the only part [\dots] which means it is singular. And since they can assert their existence
simply by, for the gluon, being contained inside what must definitely exist, because it bounds a radius in which observation is
guaranteed. And for the photons simply disappearing the moment they do. For expression versus equation, I was originally unsure if it
would end in a nice equation, so I chose expression. But I would say that it's now an equation, right? Would the amount that the Higgs field
doesn't contribute, since it is only for IM between quarks, fit? Yes, alpha does not come from it, but like established earlier, it is what
binds them together to enable witness acts.''
\item[23:43Z.] ``Typo, sorry! I meant Variable Influence Balancing Equation Else it wouldn't be a VIBE field :D'' ``I also realize that
I misread DMP as not needing the Higgs field, when instead the Higgs field IS G0 providing Existence and thus spanning an
Existence-boundary containing the Information of the IM packet.'' (The rest of that section is in the Acknowledgements.) ``For black
holes, since the only thing escaping Existence getting dragged is Information, the term should be accurate enough to coin now: Let's call
it Hawking-Penrose-Information radiation.''
\end{description}

\section{The checks}\label{app:checks}

\texttt{python3 docs/balance/checks.py} needs only numpy and prints 56 checks, all of which pass. B1: e-folds for flatness. B2: the Hubble
rate's fall to its de Sitter floor and the onset of acceleration. B3: Shor's period finding for 15, coherent and recorded. B4: ergosphere,
horizon and extractable spin energy. B5: conservation in seven decays. B6: the electron's three lengths. B7: the colour singlet. B8: the
unitarity bound without a Higgs. B9: the electroweak mixing angle and the strong coupling's growth. B10: Schwinger's term. B11: $\alpha h/c$
as mass times circumference. B12: the gluons' crossing time. B13: the Higgs lifetime and the weakness of gravity's coupling. B14: the three
couplings near $10^{14}$\,GeV. B15: the proton's electromagnetic self-energy. B16: free fall weighing only the Higgs share, against
MICROSCOPE. B17: the copying condition behind Lemma~\ref{lem:classical}. B18: Brans--Dicke's $\gamma$ and Cassini's bound behind
Proposition~\ref{pr:gamma}. B19: the Hawking temperature of a solar-mass black hole.

\begin{thebibliography}{99}\small
\bibitem{dmp} S. Hrovat, \emph{The Dark Matter Protocol}, version 0.9.9, \url{https://github.com/StevenHrovat/dmp}, commit
\texttt{4400cce} (2026).
\bibitem{aldrovandi} R. Aldrovandi, J.~G. Pereira, \emph{Teleparallel Gravity: An Introduction} (Springer, Dordrecht, 2013).
\bibitem{bekenstein82} J.~D. Bekenstein, ``Fine-structure constant: Is it really a constant?'', \emph{Phys. Rev. D} 25, 1527 (1982).
\bibitem{wz} W.~K. Wootters, W.~H. Zurek, ``A single quantum cannot be cloned'', \emph{Nature} 299, 802 (1982).
\bibitem{barnum} H. Barnum, C.~M. Caves, C.~A. Fuchs, R. Jozsa, B. Schumacher, ``Noncommuting mixed states cannot be broadcast'',
\emph{Phys. Rev. Lett.} 76, 2818 (1996).
\bibitem{ggs} T.~D. Galley, F. Giacomini, J.~H. Selby, ``A no-go theorem on the nature of the gravitational field beyond quantum
theory'', \emph{Quantum} 6, 779 (2022).
\bibitem{guth} A.~H. Guth, ``Inflationary universe: a possible solution to the horizon and flatness problems'', \emph{Phys. Rev. D} 23, 347
(1981).
\bibitem{planck18} Planck Collaboration, ``Planck 2018 results. VI. Cosmological parameters'', \emph{Astron. Astrophys.} 641, A6 (2020).
\bibitem{desi2} DESI Collaboration, ``DESI DR2 results II: measurements of baryon acoustic oscillations and cosmological constraints'',
\emph{Phys. Rev. D} 112, 083515 (2025), arXiv:2503.14738.
\bibitem{shor} P.~W. Shor, ``Polynomial-time algorithms for prime factorization and discrete logarithms on a quantum computer'',
\emph{SIAM J. Comput.} 26, 1484 (1997).
\bibitem{penrose69} R. Penrose, ``Gravitational collapse: the role of general relativity'', \emph{Riv. Nuovo Cimento} 1, 252 (1969).
\bibitem{bz} R.~D. Blandford, R.~L. Znajek, ``Electromagnetic extraction of energy from Kerr black holes'', \emph{Mon. Not. R. Astron.
Soc.} 179, 433 (1977).
\bibitem{christodoulou} D. Christodoulou, ``Reversible and irreversible transformations in black-hole physics'', \emph{Phys. Rev. Lett.}
25, 1596 (1970).
\bibitem{pdg} S. Navas et al. (Particle Data Group), ``Review of particle physics'', \emph{Phys. Rev. D} 110, 030001 (2024).
\bibitem{hensen} B. Hensen et al., ``Loophole-free Bell inequality violation using electron spins separated by 1.3 kilometres'',
\emph{Nature} 526, 682 (2015).
\bibitem{rosenfeld} W. Rosenfeld et al., ``Event-ready Bell test using entangled atoms simultaneously closing detection and locality
loopholes'', \emph{Phys. Rev. Lett.} 119, 010402 (2017).
\bibitem{storz} S. Storz et al., ``Loophole-free Bell inequality violation with superconducting circuits'', \emph{Nature} 617, 265 (2023).
\bibitem{atlas12} ATLAS Collaboration, ``Observation of a new particle in the search for the Standard Model Higgs boson with the ATLAS
detector at the LHC'', \emph{Phys. Lett. B} 716, 1 (2012).
\bibitem{cms12} CMS Collaboration, ``Observation of a new boson at a mass of 125 GeV with the CMS experiment at the LHC'', \emph{Phys.
Lett. B} 716, 30 (2012).
\bibitem{atlas22} ATLAS Collaboration, ``A detailed map of Higgs boson interactions by the ATLAS experiment ten years after the
discovery'', \emph{Nature} 607, 52 (2022).
\bibitem{cms22} CMS Collaboration, ``A portrait of the Higgs boson by the CMS experiment ten years after the discovery'', \emph{Nature}
607, 60 (2022).
\bibitem{lqt} B.~W. Lee, C. Quigg, H.~B. Thacker, ``Weak interactions at very high energies: the role of the Higgs-boson mass'',
\emph{Phys. Rev. D} 16, 1519 (1977).
\bibitem{glashow} S.~L. Glashow, ``Partial-symmetries of weak interactions'', \emph{Nucl. Phys.} 22, 579 (1961).
\bibitem{weinberg} S. Weinberg, ``A model of leptons'', \emph{Phys. Rev. Lett.} 19, 1264 (1967).
\bibitem{salam} A. Salam, ``Weak and electromagnetic interactions'', in \emph{Elementary Particle Theory}, ed. N. Svartholm (Almqvist and
Wiksell, Stockholm, 1968), 367.
\bibitem{schwinger} J. Schwinger, ``On quantum-electrodynamics and the magnetic moment of the electron'', \emph{Phys. Rev.} 73, 416
(1948).
\bibitem{fan} X. Fan, T.~G. Myers, B.~A.~D. Sukra, G. Gabrielse, ``Measurement of the electron magnetic moment'', \emph{Phys. Rev. Lett.}
130, 071801 (2023).
\bibitem{weinberg64} S. Weinberg, ``Photons and gravitons in S-matrix theory: derivation of charge conservation and equality of
gravitational and inertial mass'', \emph{Phys. Rev.} 135, B1049 (1964).
\bibitem{deser70} S. Deser, ``Self-interaction and gauge invariance'', \emph{Gen. Relativ. Gravit.} 1, 9 (1970).
\bibitem{yang} Y.-B. Yang, J. Liang, Y.-J. Bi, Y. Chen, T. Draper, K.-F. Liu, Z. Liu, ``Proton mass decomposition from the QCD energy
momentum tensor'', \emph{Phys. Rev. Lett.} 121, 212001 (2018).
\bibitem{jz} E. Joos, H.~D. Zeh, ``The emergence of classical properties through interaction with the environment'', \emph{Z. Phys. B}
59, 223 (1985).
\bibitem{zurek} W.~H. Zurek, ``Quantum Darwinism'', \emph{Nat. Phys.} 5, 181 (2009).
\bibitem{farhi} E. Farhi, J. Goldstone, S. Gutmann, J. Lapan, A. Lundgren, D. Preda, ``A quantum adiabatic evolution algorithm applied
to random instances of an NP-complete problem'', \emph{Science} 292, 472 (2001).
\bibitem{aharonov} D. Aharonov, W. van Dam, J. Kempe, Z. Landau, S. Lloyd, O. Regev, ``Adiabatic quantum computation is equivalent to
standard quantum computation'', \emph{SIAM J. Comput.} 37, 166 (2007).
\bibitem{microscope} P. Touboul et al. (MICROSCOPE Collaboration), ``MICROSCOPE mission: final results of the test of the equivalence
principle'', \emph{Phys. Rev. Lett.} 129, 121102 (2022).
\bibitem{cassini} B. Bertotti, L. Iess, P. Tortora, ``A test of general relativity using radio links with the Cassini spacecraft'',
\emph{Nature} 425, 374 (2003).
\bibitem{will} C.~M. Will, ``The confrontation between general relativity and experiment'', \emph{Living Rev. Relativ.} 17, 4 (2014).
\bibitem{bd} C. Brans, R.~H. Dicke, ``Mach's principle and a relativistic theory of gravitation'', \emph{Phys. Rev.} 124, 925 (1961).
\bibitem{higgs64} P.~W. Higgs, ``Broken symmetries and the masses of gauge bosons'', \emph{Phys. Rev. Lett.} 13, 508 (1964).
\bibitem{be64} F. Englert, R. Brout, ``Broken symmetry and the mass of gauge vector mesons'', \emph{Phys. Rev. Lett.} 13, 321 (1964).
\bibitem{hawking} S.~W. Hawking, ``Particle creation by black holes'', \emph{Commun. Math. Phys.} 43, 199 (1975).
\bibitem{oppenheim} J. Oppenheim, ``A postquantum theory of classical gravity?'', \emph{Phys. Rev. X} 13, 041040 (2023).
\bibitem{bose} S. Bose et al., ``Spin entanglement witness for quantum gravity'', \emph{Phys. Rev. Lett.} 119, 240401 (2017).
\bibitem{mv} C. Marletto, V. Vedral, ``Gravitationally induced entanglement between two massive particles is sufficient evidence of
quantum effects in gravity'', \emph{Phys. Rev. Lett.} 119, 240402 (2017).
\end{thebibliography}

\end{document}
